%!TEX root = ../thesis.tex

\chapter{Discussion: Reproducibility Across Implementations}
\label{chap:reproducibility-discussion}

No validation check in Chapter~\ref{chap:mhd-validation-results} found the
implementation failure it targets, and none measures physical accuracy, so a
matched difference clearing them is a response to the declared axis with its
mechanism left open, not an error.

\section{Relative importance of the tested axes}
\label{sec:ch6-relative-importance}

Importance is case- and criterion-specific, and a small domain average can hide
a large local one: Orszag--Tang/HLLD at $512^2$ gave
$L_\infty=5.031\times10^{-2}$ against a mean-$L_1$ of $5.554\times10^{-5}$. That
shocks and current sheets occupy few cells, so displacing one dominates
$L_\infty$, holds for HLLD but not generally: the mapped
discrepancy of Section~\ref{sec:ch5-cross-system} puts $38\%$ of HLLD's total on
$1\%$ of cells and only $6\%$ of HLL's, the diffusive clamped path spreading its
discrepancy as it spreads the features.

HLL's equality branches are unreachable by construction rather than by luck
(Section~\ref{sec:ch5-build-sensitivity}), so the branch response HLLD shows must
originate at its intermediate ties, and instrumenting them narrows it to one: of
the three, only $S_M$ ever evaluates to
exactly zero on Brio--Wu, at $62\%$ of interface evaluations, a share the
initial state's uniform zero velocity makes available and one Orszag--Tang never
reaches at all. At a matched time step HLLD also attains the higher
refinement rate on both two-dimensional cases
(Section~\ref{sec:ch4-orszag-tang}), so the sensitivity is not bought by
inaccuracy: the more accurate solver is here the more implementation-sensitive
one, through one specific decision rather than through its elaborateness.

The null CPU/GPU saved-state differences establish deterministic equivalence for
those matched configurations, but it was arranged rather than found: it needs
device contraction disabled to match the host, and restoring nvcc's default
removed it in every pair (Section~\ref{sec:ch5-hardware}). Bitwise agreement
across processor types is a property of the build configuration, not of the
hardware. Nor does the larger difference make the contracted result worse, since
a fused multiply--add rounds once where the host rounds twice; the axis is
ranked by displacement from the baseline, not by accuracy, which is why the
still more relaxed \texttt{--use\_fast\_math} build sits closer to the host.
Widening the registers changed nothing, because the semantics rather than the
instruction mix decide. These distinctions stop ``stable'' becoming shorthand
for ``reproducible'': completion, a changed saved field and an identical stored
representation are three different claims.

On the clamped HLL path at $128^2$ all three cases climb to two or three
roundings in the density domain mean and then stop, so the expectation that a
chaotic two-dimensional problem must separate faster than a one-dimensional
shock tube fails because there is a level to reach, Brio--Wu having been stopped
on the way to it (Section~\ref{sec:ch5-temporal}). Repeating that series one
factor at a time locates what the plateau belongs to
(Section~\ref{sec:ch5-saturation-axes}): not the time step, which changes the
level and not the shape, and not the arithmetic, since refining the grid doubles
the fitted exponent, replacing the clamped flux with HLLD quadruples it and
removes the plateau, and the two compound. Bounded separation is a property of
the most dissipative configuration in the matrix, which the least dissipative
path contradicts, so the reading survives only where a configuration is stated
with it.

One limit stops that becoming a claim about ideal MHD: the amplifying window
lies past the case's validation stopping time, so it is covered by completion
and physical-state checks and by no reference. The obvious competing
explanation, that those runs are simply losing divergence control, would require
the divergence residual to track the exponent, and refinement moves them
oppositely --- halving HLLD's residual while raising its exponent --- so the
growth is not the constraint decaying.

Performance is workload-dependent: in cell updates per second the CPU rate
barely moved between the two problems while the GPU rate spanned seventyfold,
which is what amortising fixed launch, transfer and output costs looks like.
Every CPU rate here is single-core, so the crossover would move with a
work-shared host. The modest precision speed-ups suit a bandwidth-bound
placement \citep{williamsEtAl2009roofline}, but with no operational intensity or
hardware counter measured that placement is inferred rather than established.

Table~\ref{tab:ch6-axis-ranking} collects the responses by axis.

\begin{table}[htbp]
\centering
\small
\begin{tabularx}{\textwidth}{@{}p{0.16\textwidth}XXX@{}}
\toprule
Axis & Domain-mean response & Local ($L_\infty$) response & Case dependence \\
\midrule
Precision (fp32/fp64) & 2.3--6.3 $u_{32}$ on coarse grids; 335 $u_{32}$ at
Orszag--Tang/HLLD $512^2$ & up to $5.031\times10^{-2}$ & strong in case and
grid \\
Build semantics (/fp:fast) & zero at fp64 scale; $3.51\,u_{32}$ for fp32
Brio--Wu & fp32 $1.311\times10^{-6}$
(HLL), $1.729\times10^{-6}$ (HLLD) & both solvers; in fp32 exceeded the
precision response for LW3 \\
Device FMA contraction (nvcc) & $3.16\,u_{32}$ for fp32 Brio--Wu & fp32
$2.265\times10^{-6}$ (Brio--Wu), $2.074\times10^{-5}$ (Orszag--Tang) & removes
CPU/GPU bitwise agreement in all four pairs \\
Branch rule ($\le$/$<$) & zero for HLL$^{\ast}$ & zero for HLL$^{\ast}$ by
construction; fp32 HLLD $2.608\times10^{-6}$ & solver-determined \\
CFL number & not measured & fp32--fp64 gap falls fourfold in step count and by
5.25 (HLL) or 2.28 (HLLD) from CFL 0.2 to 0.8 & solver-determined; $n=1$ per
setting \\
Optimisation level (/O2, /Ox) & zero & zero & none measured \\
Instruction set (/arch:SSE2, /arch:AVX2) & zero & zero, bitwise & none
measured \\
Processor type (CPU/GPU, HLL$^{\ast}$) & zero & zero, bitwise & at matched
contraction only \\
Thread count & zero & zero, bitwise & 1--8 threads; inactive in the MHD path,
a measured null on the work-shared Euler path at $4.79\times$ speed-up \\
\bottomrule
\end{tabularx}
\caption[Response by controlled axis]{\textbf{Response by controlled axis.}
Domain means are mean-relative $L_1$, local values absolute $L_\infty$, so the
two columns are not on a common scale; density only. HLL$^{\ast}$ is the
globally clamped path of Section~\ref{sec:ch2-riemann}.}
\label{tab:ch6-axis-ranking}
\end{table}

On the four axes sharing the Brio--Wu fp32 configuration exactly, all on the
HLL path since only HLL has device coverage, device contraction gives the
largest response, then relaxed host semantics, then precision, with the branch
rule at zero; optimisation level, instruction set, processor type at matched
contraction and thread count produced no measured change. The three non-zero
responses span a factor of two and all lie within the p24 Monte Carlo spread
(Section~\ref{sec:ch5-mca}), so their sequence is not a ranking, and the two
norms disagree on it: in the domain mean relaxed host semantics lead at
$3.51\,u_{32}$ against contraction's $3.16$ and precision's $2.35$, exchanging
the top two of the $L_\infty$ order. What the configuration establishes is the
division between axes that move the stored field and axes that leave it
untouched, and on HLLD the branch rule crosses into the first group with the
largest response measured here. All of this holds for these cases and grids
under one toolchain, and for density alone.

Cross-code comparison needs shared benchmarks and diagnostics
\citep{kritsukEtAl2011comparison}, making PLUTO, Athena and Athena++
\citep{mignoneEtAl2007pluto,stoneEtAl2008athena,stoneEtAl2020athenapp,
greteEtAl2021kathena} the natural independent-implementation targets, and since
implementations of one specification can diverge far beyond anything measured
here \citep{hatton1997tExperiments}, that comparison would set these effects
against a much coarser scale.

\section{Reproducibility of a published algorithm}
\label{sec:ch6-published-algorithm}

Operationally the MUSCL--Hancock/GLM update proved reliable: across two
precisions, four CFL numbers, three grids and two devices, every reported run
finished with finite, positive states. Insensitivity is another matter, since
relaxed floating-point semantics and the HLLD branch rule each moved the saved
field, and the fp32--fp64 separation grew under refinement, and in time on the
least dissipative flux, instead of shrinking. A scheme can be perfectly stable
and still fail to reproduce bit for bit across builds
\citep{mesnardBarba2017cfd}.

That leaves the question the project began from: when is fp32 defensible? Only
the case with an exact solution can answer it, and the answer is conditional on
the grid: fp32 tracked fp64 while truncation dominated and left it at $N=2048$,
beyond which refining in fp32 bought less than moving to fp64 on half the cells
(Section~\ref{sec:ch4-cp-alfven}). Precision and resolution are chosen together,
for a scheme and a problem, not once for a code.

Reproducing a published algorithm therefore requires more than naming HLL, HLLD
or MUSCL--Hancock. The record must fix an immutable source revision, a build
recipe with the effective compiler and floating-point options, a complete case
definition, the algorithmic choices that alter the arithmetic path, and the
execution platform, acceptance criteria and timing protocol actually used, bound
to the retained artefacts by hash
\citep{plesser2018reproducibility,sandveEtAl2013reproducible} as
Appendix~\ref{sec:appendix-provenance} does. Containers preserve much of the
software environment, but not hardware-dependent arithmetic.

The record supports two distinct claims. Bitwise reproducibility requires the
same stored representation under tightly matched execution; scientific
reproducibility requires the declared conclusion to survive comparison with a
justified reference and stated tolerances, and can stand where bitwise
agreement between independent implementations is infeasible, provided
unexplained differences are exposed through retained outputs. Where no
reference is available at all, a distributional criterion can take its place:
\citet{bakerEtAl2015ensemble} accept a changed configuration when its output
falls inside an ensemble generated by perturbations already judged acceptable.
